Knowledge distillation (KD) suits camera-assisted waste sorting on low-power hardware, because only the compact student is deployed.
A recent proposal adds a quantum kernel to the distillation loss, so that quantum computation is needed only during training. Testing whether such a kernel helps beyond kernel alignment itself requires matched classical controls.
We run this test on TrashNet, with RPS-25 as a low-data confirmation, using a ResNet-50 teacher, three students and one shared pipeline. Hard-label training, logit KD, relational KD, a classical RBF-kernel twin and a projected quantum kernel (PQK) are compared with three seeds per cell.
PQK is not distinguishable from logit KD. It is ahead on the mean in \KdAheadFp{} of \NPairs{} student--dataset pairs and behind on both TrashNet MobileNets, and no paired difference is significant (smallest uncorrected $p=\MinPKd$).
At its pre-specified kernel scale it beats its RBF twin in only \GtwoFixedWins{} of \NPairs{} pairs. Its seed-to-seed variance tends to be lower than the twin's, though not significantly.
All \NIntModels{} exported MobileNet students are fully integer INT8 models of about \MbMnvTwo{} and \MbMnvThree{}~MB that run in \LatMnvTwo{} and \LatMnvThree{}~ms on a Raspberry Pi~4B, and the INT8 kernels used on the device change each method's measured effect.
The quantum kernel leaves no trace in the deployed model, but our exploratory evidence does not show an accuracy benefit over classical alternatives.
